% =========================================================================
% Master Dissertation Preamble
% Biology & Zoology Specialized Academic Template
% =========================================================================

% --- Page Geometry & Thesis Layout ---
\usepackage[
  a4paper,
  top=1.0in,
  bottom=1.0in,
  left=1.25in,
  right=1.0in,
  headheight=18pt,
  headsep=10pt,
  footskip=28pt
]{geometry}

% --- Line Spacing & Paragraph Formatting ---
\usepackage{setspace}
\usepackage{titlesec}
\usepackage{placeins}
\onehalfspacing
\setlength{\parindent}{0.35in}
\setlength{\parskip}{0pt}
\setlength{\emergencystretch}{3em}
\raggedbottom
\titleformat{\chapter}[block]
  {\normalfont\bfseries\Large}
  {\thechapter.}{0.6em}{\MakeUppercase}
\titleformat{name=\chapter,numberless}[block]
  {\normalfont\bfseries\Large\centering}
  {}{0pt}{\MakeUppercase}
\titlespacing*{\chapter}{0pt}{4pt}{16pt}
\titleformat{\section}[block]
  {\normalfont\bfseries\large}
  {\thesection}{0.6em}{}
\titlespacing*{\section}{0pt}{12pt}{8pt}
\titleformat{\subsection}[block]
  {\normalfont\bfseries\normalsize}
  {\thesubsection}{0.6em}{}
\titlespacing*{\subsection}{0pt}{10pt}{6pt}

% --- Micro-typography & Fonts ---
\usepackage[utf8]{inputenc}
\usepackage[T1]{fontenc}
\usepackage{lmodern}
\usepackage{microtype}

% --- Mathematical Typesetting ---
\usepackage{amsmath}
\usepackage{amssymb}
\usepackage{amsfonts}
\usepackage{amsthm}

% --- Tables & Multi-page Catalogs ---
\usepackage{booktabs}
\usepackage{longtable}
\usepackage{tabularx}
\usepackage{array}
\usepackage{multirow}

% --- Graphics & Multi-panel Figures ---
\usepackage{graphicx}
\graphicspath{{figures/}}
\usepackage{caption}
\usepackage{subcaption}
\usepackage{float}
\captionsetup{
  font={small,stretch=1.1},
  labelfont=bf,
  format=plain,
  labelsep=period
}
\captionsetup[figure]{name={Fig.}}
\captionsetup[table]{name={Table}}
% Dissertation figures are sequential across chapters (not Fig. 3.1).
\counterwithout{figure}{chapter}

% --- Color Palette ---
\usepackage[table]{xcolor}
\definecolor{dissertationblue}{RGB}{20, 50, 120}
\definecolor{dissertationgreen}{RGB}{25, 100, 40}
\definecolor{dissertationgray}{RGB}{90, 90, 90}

% --- Phylogenetic Trees & Vector Diagrams ---
\usepackage{tikz}
\usetikzlibrary{trees,arrows.meta,positioning,shapes.geometric,calc,decorations.pathreplacing}
\usepackage{pgfplots}
\pgfplotsset{compat=1.17}
\usepackage[linguistics]{forest}

% --- Scientific Units (siunitx) ---
\usepackage{siunitx}
\sisetup{
  separate-uncertainty = true,
  multi-part-units = single,
  range-phrase = --,
  range-units = single,
  detect-all = true,
  table-alignment-mode = format,
  tight-spacing = true
}
\DeclareSIUnit\molar{M}
\DeclareSIUnit\Molar{M}
\DeclareSIUnit\unitU{U}
\DeclareSIUnit\rpm{rpm}
\DeclareSIUnit\bp{bp}
\DeclareSIUnit\Ma{Ma}
\DeclareSIUnit\Mya{Mya}

% --- Chemical Buffer Formulas & Molecular Reagents (mhchem) ---
\usepackage[version=4]{mhchem}

% --- Draft Line Numbering (Toggle for advisor review/proofreading) ---
\usepackage{lineno}
% Uncomment \linenumbers in draft mode:
% \linenumbers

% --- Bibliography & Citations (BibLaTeX) ---
% Using backend=bibtex with authoryear style ensures native, zero-dependency
% compilation with standalone Tectonic (no external Biber binary required).
\usepackage[
  backend=bibtex,
  style=authoryear,
  natbib=true,
  dashed=false,
  maxcitenames=2,
  mincitenames=1,
  maxbibnames=99,
  giveninits=true,
  uniquename=false,
  uniquelist=false,
  sorting=nyt,
  date=year,
  doi=true,
  url=true,
  isbn=false
]{biblatex}
\addbibresource{references.bib}
% Author--year list (APA-like): no ``In:'' on journal articles;
% volume(issue); clickable DOI; ragged entries so long DOIs do not stretch
\renewbibmacro{in:}{%
  \ifentrytype{article}{}{\printtext{\bibstring{in}\addspace}}}
\renewbibmacro*{volume+number+eid}{%
  \printfield{volume}%
  \printfield{number}%
  \setunit{\addcomma\addspace}%
  \printfield{eid}}
\DeclareFieldFormat[article]{number}{\mkbibparens{#1}}
\DeclareFieldFormat{doi}{\mkbibacro{DOI}\addcolon\space\href{https://doi.org/#1}{#1}}
\AtBeginBibliography{\raggedright}

% --- Hyperlinks & PDF Metadata ---
\usepackage{hyperref}
\hypersetup{
  colorlinks=true,
  linkcolor=dissertationblue,
  citecolor=dissertationgreen,
  urlcolor=dissertationblue,
  filecolor=dissertationblue,
  pdfcreator={Tectonic Typesetting System},
  pdftitle={Antibacterial Potential of Dipsacus inermis: An in vitro Study},
  pdfauthor={Bazilla Saleem (24061119001)},
  pdfsubject={M.Sc. Zoology dissertation, University of Kashmir}
}

% --- Intelligent Cross-Referencing (cleveref must load after hyperref) ---
\usepackage[capitalize,nameinlink,noabbrev]{cleveref}
\crefname{figure}{Fig.}{Figs.}
\Crefname{figure}{Fig.}{Figs.}
\crefname{table}{Table}{Tables}
\crefname{chapter}{Chapter}{Chapters}
\crefname{section}{Section}{Sections}
\crefname{appendix}{Appendix}{Appendices}

% =========================================================================
% International Code of Zoological Nomenclature (ICZN) & Life Sciences Macros
% =========================================================================

% Format taxonomic binomials in standard biological italics:
% Example: \taxa{Ranitomeya fantastica}
\newcommand{\taxa}[1]{\textit{#1}}

% Taxon with formal ICZN authority and publication year:
% Parenthesized if author originally placed taxon in another genus (ICZN Art. 51.3)
% Example: \taxonauth{Ranitomeya sirensis}{(Aichinger, 1991)}
% Example: \taxonauth{Dendrobates tinctorius}{(Cuvier, 1797)}
\newcommand{\taxonauth}[2]{\textit{#1}~#2}

% Novel species discovery designation (ICZN Recommendation 16A):
% Example: \spnov{Ranitomeya aurata}
\newcommand{\spnov}[1]{\textit{#1}~\textbf{sp.~nov.}}

% Novel genus discovery designation:
% Example: \gennov{Cryptophryne}
\newcommand{\gennov}[1]{\textit{#1}~\textbf{gen.~nov.}}

% Museum voucher holotype designation:
% Example: \holotype{QCAZ 45210}
\newcommand{\holotype}[1]{\textbf{Holotype:} \texttt{#1}}

% Museum voucher paratype designation:
% Example: \paratype{AMNH 128940}
\newcommand{\paratype}[1]{\textbf{Paratype:} \texttt{#1}}

% Type locality description:
% Example: \typelocality{Tarapoto, San Martín, Peru (06°29'S, 76°22'W, 350 m a.s.l.)}
\newcommand{\typelocality}[1]{\textbf{Type Locality:} #1}

% Gene locus formatting (italicized by convention):
% Example: \gene{COI}, \gene{16S rRNA}, \gene{RAG1}
\newcommand{\gene}[1]{\textit{#1}}

% Protein formatting (upright Roman by convention):
% Example: \protein{COX1}, \protein{RAG1}
\newcommand{\protein}[1]{\textsf{#1}}

% Laboratory protocol theorem-style environment
\theoremstyle{definition}
\newtheorem{protocol}{Protocol}[chapter]

% Running header: dissertation title + roll number
% Running footer: department + page number
\usepackage{fancyhdr}
\newcommand{\invitro}{\textit{in vitro}}
\newcommand{\invivo}{\textit{in vivo}}
% KU Zoology M.Sc. title registers (same department layout as the 2025
% Zoology dissertation used for pagination: cover in capitals, header in
% title case, certificate/declaration in sentence case). Scientific
% exceptions: binomial italic, not all-caps; \textit{in vitro} always
% lowercase italic (never all-caps or title-cased in vitro).
\newcommand{\distitle}{Antibacterial Potential of \textit{Dipsacus inermis}: An \textit{in vitro} Study}
\newcommand{\distitlecover}{ANTIBACTERIAL POTENTIAL OF\\ \textit{Dipsacus inermis}:\\ AN \textit{in vitro} STUDY}
\newcommand{\distitlecite}{Antibacterial potential of \textit{Dipsacus inermis}: an \textit{in vitro} study}
\newcommand{\disshorttitle}{\distitle}
\newcommand{\disroll}{24061119001}
\newcommand{\disfooter}{Department of Zoology, University of Kashmir, Srinagar, 190006, J\&K}
\newcommand{\applydissheaders}{%
  \fancyhf{}%
  \fancyhead[C]{%
    \scriptsize
    \makebox[\headwidth][c]{%
      \disshorttitle\hfill Enrollment No.\ \disroll
    }%
  }%
  \fancyfoot[L]{\footnotesize \disfooter}%
  \fancyfoot[R]{\thepage}%
  \renewcommand{\headrulewidth}{0.4pt}%
  \renewcommand{\footrulewidth}{0.4pt}%
}
\pagestyle{fancy}
\applydissheaders
\fancypagestyle{frontmatterstyle}{\applydissheaders}
\fancypagestyle{plain}{\applydissheaders}
